\pdfsuppressptexinfo=-1
\pdfinfoomitdate=1
\pdftrailerid{}

\documentclass{article}

\usepackage[nonatbib]{neurips_2026}
\usepackage[numbers]{natbib}

\usepackage[utf8]{inputenc}
\usepackage[T1]{fontenc}
\usepackage{hyperref}
\usepackage{url}
\usepackage{booktabs}
\usepackage{amsfonts}
\usepackage{amsmath}
\usepackage{amssymb}
\usepackage{nicefrac}
\usepackage{microtype}
\usepackage{xcolor}
\usepackage{graphicx}
\usepackage{subcaption}
\usepackage{algorithm}
\usepackage{algorithmic}
\usepackage{multirow}
\usepackage{enumitem}
\usepackage{tikz}
\usetikzlibrary{positioning, fit, arrows.meta, backgrounds, calc, shadows, shapes.geometric, shapes.misc, decorations.pathreplacing}

\newcommand{\tinker}{\textsc{Tinker}}
\newcommand{\atropos}{\textsc{Atropos}}
\newcommand{\grpo}{\textsc{Grpo}}
\newcommand{\ours}{\textsc{RL-Finetuning Bench}}

\title{RL-Finetuning Bench: An Exploratory Workshop Artifact for GRPO-Style\\
Post-Training Across Libraries, Models, and Backends}

\author{Anonymous Author(s)\\
  Affiliation withheld for blind review\\
  \texttt{anonymous@neurips.cc}}

\newcommand{\answerPartial}[1][]{\textcolor{purple}{[Partial] #1}}

\begin{document}
\sloppy
\hbadness=10000
\vbadness=10000
\hfuzz=2pt

\maketitle

\begin{abstract}
This is an \textbf{exploratory workshop artifact note}, not a benchmark paper.
We release step-level group reward traces, Zero-Variance Fraction (ZVF)
traces, a parser-script bundle, a per-run manifest, and a structured evidence
grade table for 79 GRPO-style runs across 7 RL libraries (TRL, Stable
Baselines3, CleanRL, Tianshou, PufferLib, rl\_games, d3rlpy), the closed
\tinker{} API, and 5 model families (total parameters 0.6B--671B; active
0.6B--37B for the largest mixture-of-experts checkpoints). The release lives
at an anonymous repository documented in
Section~\ref{sec:limitations}.

\textbf{Claim hierarchy (descending strength).}
\textit{(i)} ZVF~$=1\Rightarrow$~zero gradient is mechanical
(Appendix~\ref{sec:appendix:zvf-formalization}).
\textit{(ii)} On a five-seed Qwen3-8B GSM8K held-out test, post-GRPO
$83.3\%$ vs.\ base $82.0\%$ ($t{=}1.32$, $p{=}0.26$): not significant.
\textit{(iii)} The pooled ZVF--last-10 correlation
($r{=}{-}0.769$, $N{=}15$, $p{=}0.0008$) is task-confounded; within
GSM8K only ($N{=}19$), $r{=}{+}0.40$ ($p{=}0.09$).
We report ZVF as a cheap descriptive logging quantity, not as a
predictive failure detector. Frontier-scale and tool-use rows are
single-seed case studies, released for honesty and not as benchmark numbers.
The rest of this artifact note is a candid evidence-grade table and three
discussion prompts for the workshop.
\end{abstract}
 \section{Introduction}
\label{sec:intro}

\paragraph{Artifact card.}
This workshop submission is an \textbf{artifact note} accompanying a release
that contains: \textit{(i)}~step-level aggregate reward traces (mean, peak,
last-10) for every reported run; \textit{(ii)}~per-step Zero-Variance
Fraction (ZVF) and Gradient Utilization (GU) traces for every reported run;
\textit{(iii)}~a per-run manifest (Appendix~\ref{sec:run-manifest}) listing
backend, hardware, hyperparameter mapping, and evidence grade;
\textit{(iv)}~parser scripts that decode the released JSONL into a single
tidy table; and \textit{(v)}~a documented anonymous repository at
\url{https://anonymous.4open.science/r/agentic-grpo-bench-anon-2C6B} for
code, logs, and Docker images. We deliberately do \emph{not} claim a
``first comprehensive'' GRPO benchmark, a frontier ranking, or a definitive
algorithm comparison.

\paragraph{Claim hierarchy (one sentence).}
The mechanical ZVF~$\to$~zero-gradient theorem is an A-grade definitional
claim; the held-out 5-seed Qwen3-8B GSM8K control ($+1.3\,$pp,
$p{=}0.26$) is an A-grade negative result; the pooled ZVF--last-10
correlation ($r{=}{-}0.769$, $N{=}15$) is a B$^{*}$-grade descriptive
finding that is task-confounded and collapses within-task; everything
else (frontier-scale Tinker rows, tool-use, group-size sweeps) is C-grade
case study and is released for honesty rather than as benchmark fact.

\paragraph{Why this is an artifact note rather than a benchmark.}
Reinforcement-learning post-training now drives much of modern language-model
reasoning and alignment~\citep{ouyang2022training,guo2025deepseekr1nature},
but empirical claims are hard to compare:
PPO~\citep{schulman2017proximal}, GRPO~\citep{shao2024deepseekmath}, and
DPO~\citep{rafailov2024direct} are typically reported on different model
families with different framework defaults, different rollout regimes, and
different reward parsers. Rather than declare a winner, we release the
underlying traces and a candid evidence-grade table
(Table~\ref{tab:evidence-grade}) so that downstream readers can re-pool
or stratify as they see fit. The release explicitly distinguishes
multi-seed open-stack arms (5 seeds) from single-seed closed-backend
\tinker{} case studies; the workshop paper is the honest tour of what those
artifacts do and do not support. Discussion prompts for the workshop are
in Section~\ref{sec:conclusion}.
 
\section{Setup (Compact)}
\label{sec:setup}


We instrument GRPO-style training with step-level group reward statistics
across 79 runs spanning 7 RL libraries (TRL, Stable Baselines3, CleanRL,
Tianshou, PufferLib, rl\_games, d3rlpy), the closed Tinker API, and the open
Qwen3, Qwen3.5, Llama-3.x, DeepSeek-V3.1, and Nemotron-120B model families at
total-parameter scales of 0.6B--671B (active 0.6B--37B for the largest MoE
checkpoints). Headline results use Qwen3-8B with LoRA rank 32, Adam
($\beta_1{=}0.9$, $\beta_2{=}0.95$), learning rate $10^{-4}$, and 5 seeds in
$\{42,123,456,789,1024\}$. Frontier-scale and Tinker-API runs are single-seed
and reported as case studies. The full task suite, library taxonomy,
hyperparameter mapping, and per-run manifest are deferred to
Appendix~\ref{sec:run-manifest} and surrounding sections so the main text fits
the workshop page budget.

\section{Headline Result: ZVF as a Descriptive GRPO Diagnostic}
\label{sec:zvf-analysis}


\textbf{Zero-Variance Fraction (ZVF)} at step $t$ is the fraction of prompts
for which all $G$ completions receive identical rewards:
\[
  \mathrm{ZVF}_t = \frac{1}{|\mathcal{P}_t|}\sum_{p \in \mathcal{P}_t}
  \mathbf{1}\bigl[\operatorname{Var}_{c \sim \mathcal{C}(p)}\,r(c) = 0\bigr].
\]
When $\mathrm{ZVF}_t = 1$, every prompt contributes zero gradient.
\textbf{Gradient Utilization} is $\mathrm{GU}_t = 1 - \mathrm{ZVF}_t$.

Across $N{=}15$ experiments spanning 5 model families and 0.6B--671B total
parameters (0.6B--37B active for the largest mixture-of-experts checkpoints),
mean ZVF \emph{co-occurs} with final training performance: Pearson
$r{=}-0.769$, $p{=}0.0008$ (Spearman $\rho{=}-0.784$, $p{=}0.0005$).
\textbf{The pooled correlation is confounded by task type.} Tool-use
experiments reach $\mathrm{ZVF}=100\%$ at $\sim$0\% reward, while GSM8K
experiments average $\mathrm{ZVF}=8.5\%$ at much higher reward. Stratifying
by task collapses the apparent relationship within each subgroup; we
therefore read the pooled $r=-0.769$ as ``tool-use failed and GSM8K did not''
rather than as a clean diagnostic--performance link.
Stratifying within GSM8K ($N{=}19$ Tinker GSM8K runs), the apparent monotone
relationship does not survive: Pearson $r{=}+0.40$ ($p{=}0.09$); restricted to
the seven model-varying GSM8K runs, $r{=}+0.13$ ($p{=}0.79$). The pooled
correlation reflects task identity, not within-task ZVF dynamics.
We treat ZVF as a cheap descriptive logging quantity, not as a standalone
causal predictor.

\begin{figure}[t]
\centering
\includegraphics[width=\linewidth]{figures/v2/workshop_killer}
\caption{\textbf{ZVF as a descriptive (not predictive) GRPO diagnostic.}
\textit{(a)}~Per-step ZVF heatmap over the 15 cross-stack runs; red:
$\mathrm{ZVF}{=}1$, blue: $\mathrm{ZVF}{=}0$, grey: no data. Tool-use
rows (top) form the persistent red band that drives the pooled
$r{=}{-}0.769$.
\textit{(b)}~Tool-use sentinel: $\mathrm{ZVF}{=}1$, $\mathrm{GU}{=}0$,
reward$\equiv 0$ at every step.
\textit{(c)}~Residualized within-GSM8K: per-group-demeaned ZVF vs.\
last-10; within-task slope $r_{\mathrm{within}}\approx 0$, in sharp
contrast to pooled $r{=}{-}0.769$.}
\label{fig:workshop-killer}
\end{figure}

\begin{table}[htbp]
\centering\small
\setlength{\tabcolsep}{5pt}
\caption{Statistical ablation of the ZVF--last-10 correlation under three
sample definitions. The sign of the correlation \emph{flips} once task is
held fixed, and the within-task correlations do not survive a 5\%
threshold.}
\label{tab:zvf-stat-ablation}
\begin{tabular}{@{}l r r r l@{}}
\toprule
\textbf{Sample} & \textbf{$N$} & \textbf{Pearson $r$} & \textbf{$p$} & \textbf{Note} \\
\midrule
Pooled, cross-stack          & 15 & $-0.769$ & $0.0008$ & B$^{*}$, task-confounded \\
Within-GSM8K (all)           & 19 & $+0.40$  & $0.09$   & marginal, sign flips \\
Within-GSM8K (model-varying) &  7 & $+0.13$  & $0.79$   & not significant \\
\bottomrule
\end{tabular}
\end{table}

\noindent
We do not report bootstrap CIs or permutation $p$-values for the
within-GSM8K rows here; the full bootstrap analysis is in the ZVF and
D\&B variants of this work and is out of scope for the workshop note.

\section{Compact Main Results and Evidence Grade}
\label{sec:main_results}
\label{sec:results}


Table~\ref{tab:headline} reports the three kinds of evidence in the release
separately: multi-seed open-stack runs, single-seed Tinker case studies, and
known failure probes. We deliberately do not pool them into one ranking.

\begin{table}[htbp]
\centering\small
\setlength{\tabcolsep}{4pt}
\caption{Workshop headline table, segregated by evidence grade. Open-stack
arms have 5-seed coverage; Tinker arms are single-seed (cost-constrained,
closed backend) and are reported as case studies.}
\label{tab:headline}
\begin{tabular}{@{}l l l r r r@{}}
\toprule
\textbf{Algo / stack} & \textbf{Model} & \textbf{Task} & \textbf{Seeds} & \textbf{Peak} & \textbf{Last-10} \\
\midrule
\multicolumn{6}{l}{\textit{Multi-seed open-stack (workshop headline)}} \\
\midrule
TRL-GRPO (Modal L4) & Qwen2.5-0.5B & GSM8K & 5 & 0.747$_{[0.672,0.782]}$ & 0.734 $\pm$ 0.07 \\
\midrule
\multicolumn{6}{l}{\textit{Single-seed Tinker case studies (\dag closed backend)}} \\
\midrule
GRPO\,\dag        & Qwen3-8B            & GSM8K     & 1 & 62.5\% & 34.4\% \\
GRPO\,\dag        & Llama-3.1-8B-Inst   & GSM8K     & 1 & 100\%  & 84.4\% \\
GRPO\,\dag        & DeepSeek-V3.1       & GSM8K     & 1 & 100\%  & 85.0\% \\
GRPO\,\dag        & Qwen3-235B-A22B     & GSM8K     & 1 & 100\%  & 100\%  \\
GRPO\,\dag        & Nemotron-120B       & GSM8K     & 1 & 87.5\% & 16.2\% \\
\midrule
\multicolumn{6}{l}{\textit{Failure probes (released for honesty, not as benchmark numbers)}} \\
\midrule
GRPO\,\dag        & Qwen3-32B           & Tool-use  & 1 & 0\%    & 0\%    \\
GRPO\,\dag        & Llama-3.1-8B-Inst   & Tool-use  & 1 & 0\%    & 0\%    \\
\bottomrule
\end{tabular}
\end{table}

\paragraph{Held-out negative control (workshop headline).}
On a separate five-seed Qwen3-8B GSM8K base-vs-GRPO comparison, the post-GRPO
mean is $83.3\%$ versus $82.0\%$ for the base model under the same greedy
protocol. The $+1.3$pp delta is not statistically significant
($t{=}1.32$, $p{=}0.26$). This is the workshop's clearest negative result:
strong online training reward does not by itself imply held-out gain.

\paragraph{Group size: brief.}
We ran a single-seed exploratory $G$-sweep on Qwen3-8B/GSM8K whose results
are non-monotonic across $G$; see
Appendix~\ref{sec:appendix:group-size-reconcile} for the token-budget
reanalysis. \textbf{We do not headline a winning $G$.} A 3-seed
token-budget reanalysis at $T{\ge}16$M tokens places a fitted apex near
$G{\approx}32$ for one stack and objective, but the apex shifts with $T$ and
does not conflict with smaller-$G$ optima under different cost objectives; we
deliberately avoid promoting a group-size law. A subsequent two-seed matched-budget
panel (2{,}560 rollouts per arm; $G{=}2\times160$ vs.\ $G{=}16\times20$
steps) adds one trajectory-level datum: the $G{=}2$ arms drive train
reward to $\approx0.9$--$1.0$ on the sampled pool, then end in a sustained
all-\emph{correct} zero-variance regime ($\mathrm{ZVF}\approx0.75$--$1.0$),
while $G{=}16$ ends mid-learning ($\approx0.3$--$0.5$) with
$\mathrm{ZVF}\le0.25$ --- small $G$ trades early
step-count for endgame signal exhaustion.

\begin{table}[htbp]
\centering\small
\setlength{\tabcolsep}{4pt}
\caption{Evidence grade of each main-text claim in the workshop variant.
\textbf{A}~=~multi-seed open-stack or definitional;
\textbf{B}~=~multi-seed or pooled with a known caveat;
\textbf{B$^{*}$}~=~significant in pooled sample but task-confounded;
\textbf{C}~=~single-seed or closed-backend case study.}
\label{tab:evidence-grade}
\begin{tabular}{@{}p{0.50\linewidth} l p{0.32\linewidth}@{}}
\toprule
\textbf{Claim} & \textbf{Grade} & \textbf{Notes} \\
\midrule
ZVF $\to$ zero gradient is mechanical & A (def.) & App.~\ref{sec:appendix:zvf-formalization}. \\
ZVF--last-10 $r{=}-0.769$ ($N{=}15$, pooled) & B$^{*}$ & Confounded by task type. \\
Held-out Qwen3-8B GSM8K $83.3\%$ vs.\ $82.0\%$ & A & 5 seeds; $p{=}0.26$. \\
Group-size sweep is non-monotonic; no headline winner & C & single-seed per $G$. \\
Matched-budget $G{=}2$ endgame ZVF wall & B & 2 seeds/arm; 512-tok cap. \\
Token-budget fitted apex near $G{\approx}32$ at $T{\ge}16$M & B &
3-seed appendix reanalysis; not a universal optimum. \\
Frontier-scale Tinker case studies & C & single-seed, closed backend. \\
Tool-use $0\%$ reward across runs & B & task-design failure, not capability. \\
\bottomrule
\end{tabular}
\end{table}

\section{Reproducibility and Limitations}
\label{sec:limitations}


We release pinned Docker images, deterministic seed management, and
\texttt{rliable}-based analysis scripts. Anonymous code, logs, and
checkpoints are at
\url{https://anonymous.4open.science/r/agentic-grpo-bench-anon-2C6B}.
The release distinguishes three artifact tiers explicitly:
(i)~\textbf{precomputed ZVF / GU traces}, released for every reported run;
(ii)~\textbf{step-level aggregate logs} (mean reward, peak reward,
last-10), released for every reported run;
(iii)~\textbf{raw per-group rollout JSONL} (per-prompt, per-completion
reward arrays), released for a documented subset of stacks, with the full
subset enumerated in the run manifest (Appendix~\ref{sec:run-manifest})
and the parser contract released with the anonymous artifact.

The release has known gaps that the workshop framing makes explicit. (i)~Most
Tinker experiments are single-seed and depend on a closed-source backend; we
report them descriptively. (ii)~Empirical evidence concentrates on
short-horizon GSM8K with binary verifiable rewards; tool-use, code, and
non-math coverage are sparser and yield negative or descriptive evidence
only. (iii)~Train-set reward in the Tinker case studies is \emph{descriptive
instrumentation}, not an inferential signal: our primary inferential claims
are limited to the mechanical ZVF~$\to$~zero-gradient theorem
(Appendix~\ref{sec:appendix:zvf-formalization}) and the held-out 5-seed Qwen3-8B
GSM8K test-set control ($83.3\%$ vs.\ $82.0\%$, $t{=}1.32$, $p{=}0.26$).
Tinker-derived stratified correlations (e.g., the within-GSM8K ZVF
re-analysis) are reported as artifact analyses with closed-backend caveats,
not as benchmark statistics. (iv)~Cross-library comparisons are confounded by
checkpoint, hardware, and default-config differences; we report them as
stack-level observations, not as algorithm-level statements. (v)~Direct KL
tracking was blocked by a PyTorch gradient-graph error; the corrected
implementation is in the artifact but is not validated against ZVF in the
current paper. (vi)~Verifiable rewards reduce reward-model subjectivity but
do not eliminate reward hacking, format gaming, or train-set reward
inflation; we do not claim they do.

\section{Discussion and Concrete Workshop Questions}
\label{sec:conclusion}

The workshop release is deliberately exploratory. Its strongest evidence is
short-horizon GSM8K dynamics; non-math coverage is a stress case, not a
result. We treat ZVF as a cheap descriptive diagnostic, not a predictive
failure detector, and we emphasize that the Qwen3-8B held-out control gives
a small, non-significant GRPO improvement over the base.

\paragraph{Five concrete questions for the workshop.}
We invite attendees to engage with five specific, testable questions raised
by this release rather than open-ended philosophizing:
\begin{enumerate}[leftmargin=1.4em,topsep=2pt,itemsep=2pt]
  \item Does step-level ZVF predict reward change \emph{within} a single
        run, controlling for run fixed effects, on stacks beyond Tinker?
        Our residualized within-GSM8K scatter
        (Figure~\ref{fig:workshop-killer}c) finds $r_{\mathrm{within}}$
        close to zero; is this a backend artifact or a general property?
  \item Under a fixed rollout-token budget, where does group size trade early
        step count for endgame contrast? Our two-seed $G{=}2\times160$ versus
        $G{=}16\times20$ panel ends with the small-$G$ arms at an all-correct
        ZVF wall while $G{=}16$ remains mid-learning; a separate 3-seed
        reanalysis places a fitted apex near $G{\approx}32$ for one stack and
        budget. Neither identifies a universal optimum. The needed experiment
        is a directly trained, multi-seed, token-matched sweep on open stacks.
  \item Are there tasks where high-ZVF cells \emph{do} predict downstream
        collapse, beyond GSM8K? Tool-use is a degenerate $\mathrm{ZVF}{=}1$
        case from step 1; the genuinely interesting test is on
        intermediate-density tasks (MATH-500 at low temperature, code
        with partial-credit rewards).
  \item What minimum disclosure makes a managed-API result count as
        benchmark evidence? Our Tinker rows are reproducible only at the
        level of API call scripts; should the community require GRPO
        loss / reward normalization / hardware metadata for a managed-API
        row to count as anything beyond a case study?
  \item Should ZVF-style diagnostic claims require raw per-group rollout
        JSONL by default? Our release bundles step-level aggregates for
        every run but raw JSONL only for a documented subset; we invite
        the workshop to align on a minimum raw-artifact bar.
\end{enumerate}
 


\section{Ethics, Limitations, and Broader Impact}
\label{sec:impact_standalone}

This statement is written to satisfy the NeurIPS Code of Ethics and Broader
Impact requirements. It complements the shorter Limitations and Broader
Impact subsections embedded in Section~\ref{sec:limitations} by consolidating
in one place a full dual-use analysis, itemised compute accounting, carbon
footprint estimation, data provenance disclosures, a candid acknowledgment
of the closed-source training backend on which a fraction of our headline
numbers depends, and a list of methodological limits we are aware of but
did not have resources to close within the submission window.

\subsection{Dual-Use Analysis: Misuse Risks of Reasoning RL}
\label{sec:dualuse}


\paragraph{Threat model.}
\ours{} releases (i) training scripts that apply GRPO to the GSM8K
\cite{cobbe2021training} mathematical reasoning benchmark and to the
Salesforce xLAM-Function-Calling-60k \cite{liu2024xlam} tool-use corpus,
(ii) a set of LoRA adapters published on the HuggingFace Hub
(\texttt{anonymous/tinker-rl-bench-*}), and (iii) diagnostic code for
Zero-Variance Fraction, reward-stability, and length-bias analysis. We
analyse four concrete misuse pathways these artefacts could, in principle,
enable, and we describe the existing guardrails and the residual risk we
judge acceptable for publication.

\paragraph{M1. Weaponisation of mathematical reasoning.}
The most capability-relevant artefact we release is GRPO training code that
improves grade-school arithmetic and multi-step reasoning. A malicious
operator could attempt to extend this pipeline to tasks of greater
dual-use concern, e.g., chemistry olympiad problems, cryptographic key
recovery, or financial fraud. We judge the \emph{marginal} uplift provided
by our release to be small for three reasons. First, GRPO at our training
scale does not create new capabilities; the base-model control for GSM8K
(Section~\ref{sec:results}, 82.0\%) shows that the bulk of held-out
accuracy is attributable to Qwen3-8B's pretraining. Our full GRPO run only
adds \(+1.3\) percentage points ($p{=}0.26$, not statistically significant).
Second, the public literature already contains GRPO implementations
\cite{shao2024deepseekmath, openrlhf2024, verl2024, trl2020} that are at
least as capable. Third, any reasoning uplift is bounded by the rewarder:
GSM8K rewards use exact-match against a human-written gold answer, which
does not generalise to the domains listed above without a new reward
signal, which itself requires expertise and labelled data.

\paragraph{M2. Reward hacking and the long-term alignment risk.}
GRPO, like all policy-gradient RL from a proxy reward, is vulnerable to
reward hacking: the model learns to exploit idiosyncrasies of the reward
function rather than solve the intended task \cite{skalse2022specificationgaming,
krakovna2020specification}. We observed verbosity drift in some
of our own GRPO runs, which is a concrete demonstration of this failure
mode in our own data.
Users who apply our scripts to under-specified reward functions in
production settings (e.g., a custom ``helpfulness'' classifier) may observe
models that appear to improve on held-out metrics while silently
optimising against unintended proxies. We include the Zero-Variance
Fraction and reward-stability diagnostics precisely so that downstream
users can detect such drift.

\paragraph{M3. Circumventing safety training via distillation or tool-use.}
The tool-use track trains LoRA adapters that teach a base model to emit
schema-valid JSON function calls \cite{liu2024xlam}. A motivated adversary
could in principle use the same pipeline to teach a base model to emit
jailbreak prompts as ``tool calls'' or to invoke an adversarial external
tool. We mitigate this by (a) releasing only adapters, not merged weights,
so that the safety-trained instruct checkpoint must be applied separately
and any instruct-layer guardrails remain active; (b) documenting in model
cards that the tool-use adapters are trained on a narrow schema (five
synthetic tools) and have not been safety-evaluated for open-ended function
calling; and (c) not releasing any function-calling harness that would
actually execute emitted calls.

\paragraph{M4. Compute-efficiency externalities.}
Our work argues that critic-free RL with budget-dependent group-size
selection (the optimal $G$ depends on total token budget; see
Appendix~\ref{sec:appendix:group-size-reconcile}) can be run on modest
budgets. Improved training efficiency is
itself dual-use: it lowers the cost of running adversarial fine-tuning at
scale. We judge this externality to be modest relative to the already-low
cost of fine-tuning a 0.6B--8B model with commodity LoRA recipes
\cite{hu2022lora}, but we flag it explicitly.

\paragraph{Residual risk.}
After weighing M1--M4 we judge that the reproducibility, calibration, and
pedagogical benefits of releasing \ours{} outweigh the residual misuse
risk. We adopt the standard mitigation of releasing adapters (not merged
weights), documenting intended use and out-of-scope use in model cards,
and publishing alongside the diagnostics a researcher would need to detect
reward hacking.

\subsection{Compute Cost Accounting}
\label{sec:compute_accounting}


Table~\ref{tab:compute_spend} itemises the financial cost of every
platform used in the project. Costs are real spend for the submitting
authors; no in-kind credits are omitted. Dollar figures are in USD.

\begin{table}[ht]
\centering
\small
\caption{Itemised compute spend across platforms. Totals include both
successful and failed runs, because failed-run spend is real and should
not be hidden from reviewers \cite{henderson2018deep}.}
\label{tab:compute_spend}
\begin{tabular}{llrr}
\toprule
\textbf{Platform} & \textbf{Use} & \textbf{Runs} & \textbf{Cost (USD)} \\
\midrule
Tinker SDK v0.16.1 & GRPO on 4B/8B (original suite)        & 25 & \$40--55 \\
Tinker SDK v0.16.1 & 10x Structural Ceiling ablation       & 32 & \$65 \\
Tinker SDK v0.16.1 & Frontier Suite (frontier/MoE)      & 13 & \$25--30 \\
Modal H100         & PPO baselines (Qwen3-8B, Llama-3.1-8B)& 2  & \$12--18 \\
Modal H100         & KL-tracking run (failed)              & 1  & \$2--4   \\
Google Colab Pro (T4) & 0.5B--3B QLoRA SFT+GRPO            & --- & \$10/person \\
NVIDIA L4 (TRL baseline)& Qwen2.5-0.5B GSM8K, 5 seeds      & 5  & \$3--5   \\
HuggingFace Hub    & Model + adapter hosting               & --- & \$0 \\
Weights \& Biases  & Experiment tracking (academic tier)   & --- & \$0 \\
\midrule
\multicolumn{3}{l}{\textbf{Total authors' out-of-pocket spend}} & \textbf{\$130--140} \\
\bottomrule
\end{tabular}
\end{table}

\paragraph{Hidden costs.} Beyond direct platform spend, the project
consumed human labour (approximately six student-FTE-months, unpaid) and
infrastructure generously subsidised by an academic research group with
shared access to one A100 node used for TRL baselines and
pre-submission sanity checks). No author received commercial compensation
from Thinking Machines, Modal, or any other vendor mentioned.

\paragraph{Failed-run accounting.} Of the Tinker spend, we estimate
\textasciitilde\$30--35 was consumed by runs that ultimately failed
(JWT expiry, API stalls for GPT-OSS-20b and Kimi-K2 before re-run,
KL-tracking gradient bug). We disclose this because reviewers often ask
whether reported total spend reflects only clean, successful runs; it does
not. Excluding failed spend would understate the true resource cost of
reproducing our work by roughly 25\%.

\subsection{Carbon Footprint Estimation}
\label{sec:carbon}


Table~\ref{tab:carbon2} computes an energy and CO$_2$-equivalent estimate
for each platform. We follow the methodology of
\citet{patterson2024empirical} and \citet{strubell2019energy}: for each
GPU-hour we multiply device TDP by wall-clock GPU-hours and the data-centre
Power Usage Effectiveness (PUE), then multiply by an anonymized regional
average grid intensity for the cloud zone in which the hardware is
believed to be located.

\paragraph{Assumed constants.}
\begin{itemize}
  \item NVIDIA H100 SXM5 TDP: 700\,W.
  \item NVIDIA A100 80GB TDP: 400\,W.
  \item NVIDIA L4 TDP: 72\,W.
  \item NVIDIA T4 TDP: 70\,W.
  \item Data-centre PUE: 1.1 (conservative; typical hyperscale PUE is
        1.10--1.15 \cite{google2023pue, aws2023pue}).
  \item US average grid intensity (EPA eGRID 2023): 0.367\,kg\,CO$_2$-eq / kWh
        \cite{epa2023egrid}.
  \item Author-region average grid intensity (national authority report,
        2023; specific country and source withheld for anonymity):
        within the higher-carbon-intensity grid bands typical of
        emerging-market national grids; the exact value used is the
        midpoint of an anonymized regional average grid intensity range
        and is withheld here.
  \item We use the US average for Modal H100 (an anonymized regional
        cloud zone as configured) and Tinker (location undisclosed;
        assumed US), and an anonymized regional average grid intensity
        for the Colab Pro T4 instances used by the authors. The exact
        regional intensity is withheld for anonymity; the reported
        project-total CO$_2$-eq number uses the midpoint of the
        applicable range.
\end{itemize}

\paragraph{GPU-hour estimates.}
Tinker GPU-hour counts are not disclosed by the platform. We estimate them
from wall-clock run duration and assume a single H100-class accelerator per
job, which is consistent with Tinker's published architecture for the
model sizes we trained. Modal GPU-hours are measured directly from Modal's
billing dashboard.

\begin{table}[ht]
\centering
\small
\caption{Energy and carbon footprint estimates. Tinker hardware is not
publicly disclosed; we assume 1$\times$ H100 equivalent per run. Failed /
stalled runs are included. Totals rounded to one significant figure.}
\label{tab:carbon2}
\begin{tabular}{lrrrrr}
\toprule
\textbf{Platform} & \textbf{GPU-h} & \textbf{TDP\,(W)} & \textbf{PUE} &
\textbf{Energy\,(kWh)} & \textbf{CO$_2$\,(kg)} \\
\midrule
Tinker (assumed H100)   & \textasciitilde 950 & 700 & 1.1 & \textasciitilde 732 & \textasciitilde 269 \\
Modal H100 (anonymized zone)    & \textasciitilde 18  & 700 & 1.1 & \textasciitilde 14  & \textasciitilde 5.1 \\
Institutional A100 (in-kind)      & \textasciitilde 60  & 400 & 1.1 & \textasciitilde 26  & \textasciitilde 19   \\
Colab Pro T4 (regional)    & \textasciitilde 40 & 70 & 1.2 & \textasciitilde 3.4 & \textasciitilde 2.4  \\
NVIDIA L4 (TRL baseline)   & \textasciitilde 10 & 72 & 1.1 & \textasciitilde 0.8 & \textasciitilde 0.3  \\
\midrule
\textbf{Project total (best estimate)} & & & & \textasciitilde 776 & \textasciitilde 296 \\
\bottomrule
\end{tabular}
\end{table}

\paragraph{Interpretation and uncertainty.}
Our best estimate of \textasciitilde 296\,kg\,CO$_2$-eq for the entire
project is comparable to a single round-trip domestic economy flight in
the same region (\textasciitilde 300\,kg). The dominant uncertainty is
the Tinker GPU-hour count: because Tinker does not expose hardware telemetry,
a factor-of-two error in GPU-hours or TDP is plausible. Under a pessimistic
assumption (2$\times$ GPU-hours, H100 running near 100\% TDP continuously),
the Tinker contribution could be as high as \textasciitilde 540\,kg CO$_2$-eq.
Under an optimistic assumption (Tinker backends actually use more
energy-efficient accelerators than H100, 50\% average utilisation) the
contribution could be as low as \textasciitilde 130\,kg. We report the
central estimate in the main table and caution readers that it should not
be interpreted as precise. \emph{All} numbers should be treated as
order-of-magnitude estimates with documented sensitivity to Tinker
hardware assumptions; subsequent users can additionally skip the failed
runs and the ablation sweeps to lower the cost of reproducing our work.

\paragraph{Offset and mitigation.}
We did not purchase carbon offsets. Instead, we have (a) released all
trained checkpoints on HuggingFace Hub so that downstream users do not
need to re-run our experiments; (b) released step-level CSV logs so that
learning curves can be inspected without re-training; and (c) documented
in Section~\ref{sec:compute_accounting} which runs were wasteful, so that
replicators can skip them. We regard reproducibility-by-artefact as a
more durable mitigation than offsets.

\subsection{Data Provenance}
\label{sec:provenance}


\ours{} uses only publicly released research datasets. No private data,
personally identifiable information, or licensed proprietary content is
used. No human annotators were employed during this work. We document
each dataset below.

\paragraph{GSM8K \cite{cobbe2021training}.}
8{,}500 grade-school math word problems (7{,}473 train / 1{,}319 test)
authored by human writers contracted by OpenAI. Released under the
\textbf{MIT License} via \url{https://github.com/openai/grade-school-math}.
We use the standard train/test splits without modification. The canonical
citation is \citet{cobbe2021training}. Known limitations of GSM8K include
gender-neutral but culturally US-centric word problems (names, currencies,
sports) and a moderate rate of human-labelling errors estimated at
\textasciitilde 2\% by \cite{zhang2024gsm1k}. Our held-out evaluation
respects the test/train split.

\paragraph{Salesforce xLAM-Function-Calling-60k \cite{liu2024xlam}.}
60{,}000 function-calling examples generated by Salesforce Research as
training data for the xLAM agentic model family. Released under the
\textbf{CC BY 4.0} license via \url{https://huggingface.co/datasets/Salesforce/xlam-function-calling-60k}.
We use the public release as-is; specifically, we use the first 35 prompts
$\times$ 10 rollouts in the 10x Structural Ceiling tool-use track and the
full 60k split for the xlam-60k real-data run in Section~\ref{sec:results}.
Known limitations: xLAM schemas are synthetic and skew toward cleanly-typed
arguments; the dataset under-represents ambiguous, error-handling, and
multi-turn tool-calling. Attribution to Salesforce is required by the
license and is provided in Section~\ref{sec:acknowledgments} and in the
xLAM-derived model cards.

\paragraph{HumanEval \cite{chen2021codex}.}
164 Python programming problems with unit tests, released by OpenAI under
the \textbf{MIT License}. Used unmodified for pass@$k$ evaluation. Known
limitations: small scale, English-language docstrings, single-language
coverage; documented contamination concerns against frontier pretraining
corpora \cite{riddell2024contamination}.

\paragraph{NuminaMath \cite{li2024numinamath}.}
Used by an anonymous collaborator for the multi-stage
GSM8K+NuminaMath pipeline. Released under \textbf{Apache 2.0}. We use a
downstream checkpoint reported by that collaborator; our repository contains only their
scripts and the Apache-licensed derivative weights.

\paragraph{Open-Platypus \cite{lee2023platypus}.}
3{,}000 SFT examples used by an anonymous collaborator for Qwen3-8B code
generation warm-up. Open-Platypus is a curated subset released under
\textbf{CC BY-NC 4.0}; the non-commercial restriction is respected in our
release, which is academic-use only.

\paragraph{Synthetic tool-use corpus (authored).}
We additionally generated a small five-tool synthetic corpus (calculator,
web-search stub, calendar, file-read, email-send) used for the tool-use
case studies cited in Section~\ref{sec:zvf-analysis} (the high-ZVF /
zero-reward tool-use cluster that drives the pooled task confound).
The corpus is authored by the paper's authors from
publicly documented APIs, contains no third-party content, and is
released under the \textbf{MIT License} in our repository under
\texttt{data/synthetic\_tools/}. Generation prompts are included for
auditability. No user data, scraped web content, or proprietary API
traces are present.

\paragraph{No web scraping, no human subjects.}
We did not scrape any website. We did not run any study with human
participants. No IRB review was therefore required. No data that could
reasonably be considered to contain PII, copyrighted content, or sensitive
personal attributes was used at any stage of training, evaluation, or
reward modelling.

\subsection{Closed-Source Tinker Acknowledgment}
\label{sec:tinker_closed}


A central tension in \ours{} is that a substantial fraction of our
headline numbers come from a closed-source commercial platform while our
paper simultaneously advocates for reproducibility and platform
independence. We do not resolve this tension by hiding it; we describe
it precisely.

\paragraph{What Tinker is and is not.}
Tinker \cite{thinkingmachines2024tinker} is a managed LLM fine-tuning and
inference service provided by Thinking Machines, Inc. It exposes a Python
SDK (\texttt{tinker==0.16.1}) that accepts custom loss functions
(\texttt{forward\_backward\_custom}), a limited set of optimisers, and
standard LoRA hyperparameters. It does not expose (i) the exact server-side
GRPO loss implementation, (ii) reward normalisation or baseline subtraction
scheme, (iii) minibatch construction or gradient-accumulation strategy,
(iv) hardware configuration (GPU type, inter-node bandwidth), or (v)
system-level telemetry (energy, throughput, queueing).

\paragraph{What this means for our claims.}
Tinker results in this paper measure the \emph{platform's} implementation
of GRPO, not an abstract specification of the algorithm. Cross-stack reward
gaps between Tinker GRPO and an open-source GRPO implementation cannot be
fully attributed from our data: we are able to rule out a handful of
candidate explanations (seed variance, model-size confound) but not the
implementation itself. We therefore draw \emph{quantitative} conclusions only
from the open-source side (TRL, veRL on Modal H100, OpenRLHF) where every
hyperparameter is auditable, and use Tinker results primarily as
\emph{descriptive case studies} of what a carefully engineered production
stack can achieve for critic-free RL at our scales.

\paragraph{Reproducibility commitments we can make.}
\begin{enumerate}
  \item All Tinker experiment scripts, configuration files, JSON step
        logs, and per-run W\&B projects are archived in the repository.
        Researchers with Tinker access can attempt replication with our
        exact configurations.
  \item Every figure and every summary statistic derived solely from
        Tinker data is marked with ``\dag'' and a footnote stating that
        independent replication requires Tinker API access.
  \item Primary inferential claims are limited to (i) the mechanical
        zero-gradient theorem (Appendix~\ref{sec:appendix:zvf-formalization})
        and (ii) the held-out 5-seed Qwen3-8B GSM8K control reported in
        Section~\ref{sec:limitations}. The Tinker-derived stratified ZVF
        correlations (Section~\ref{sec:zvf-analysis}) are reported as
        artifact analyses with explicit closed-backend caveats and
        single-seed-per-cell limitations, not as platform-independent
        statistical conclusions.
  \item We commit to re-running any Tinker-only experiment on an open
        backend if an equivalent open-source service becomes available,
        and to issuing a revised version of this paper if any Tinker
        result is ever shown to be due to an undisclosed
        implementation choice rather than the algorithm.
\end{enumerate}

\paragraph{Key rotation and credential hygiene.}
Our Tinker API key (prefix \texttt{tml-...}) is held in a repository
\texttt{.env.example} placeholder only. The real key is stored in the
authors' password manager and rotated whenever a failure mode suggests
possible exfiltration. No Tinker key is committed to git history in the
anonymous source repository
(\url{https://anonymous.4open.science/r/agentic-grpo-bench-anon-2C6B}).

\subsection{Known Methodological Limits}
\label{sec:methodlimits}


In addition to the infrastructure failures and platform-specific caveats
enumerated in Section~\ref{sec:limitations}, we flag the following
methodological limits.

\paragraph{Short training horizons.}
All Tinker GRPO runs were capped at 30--50 gradient steps. This is a
deliberate cost-control choice and means our Tinker numbers are
\emph{early-training snapshots}. We cannot rule out that longer training
would change the qualitative picture (e.g., the 82\%$\to$83.3\% held-out
gain might widen, or might collapse to reward hacking).

\paragraph{Single-seed Tinker experiments.}
Each Tinker configuration was run once, not 3--5 times as best practice
prescribes \cite{henderson2018deep}. Tinker results therefore lack
variance estimates and do not support significance testing; we report
them descriptively. Only the 5-seed TRL baseline and the 5-seed held-out
GSM8K evaluation carry proper confidence intervals.

\paragraph{Train-set reward as primary Tinker metric.}
Tinker runs primarily report reward on training prompts. Only the GSM8K
track was followed up with a separate 200-example held-out evaluation
per seed. Other Tinker tracks (tool-use, xLAM) remain train-set-only and
we cannot distinguish memorisation from generalisation for those
settings.

\paragraph{LoRA only; no full fine-tuning.}
All experiments use LoRA \cite{hu2022lora} with ranks 8--64. We have
\emph{not} tested whether full fine-tuning would re-order the
library/algorithm comparisons. The LoRA constraint is practical (cost) but
it means our claims about PPO vs.\ GRPO and TRL vs.\ Tinker are
LoRA-conditional.

\paragraph{Benchmark coverage is narrow.}
We evaluate four domains: GSM8K, MATH-500 (exploratory), HumanEval
(subset), and synthetic/xLAM tool-use. We do not evaluate on MT-Bench,
ArenaHard, safety benchmarks (HarmBench, ToxicChat), or truthfulness
benchmarks (TruthfulQA). Any claim about ``reasoning'' in the paper is
strictly a claim about the covered benchmarks.

\paragraph{No human preference data.}
We use verifiable rewards only (exact-match for math, unit-test for code,
schema-validity for tool-use). We do not study RLHF with human preference
data; findings should not be extrapolated to reward-model--based alignment
without further work.

\paragraph{Closed-source implementation opacity (reiterated).}
Comparisons between Tinker GRPO and TRL GRPO are confounded by the
closed-source nature of Tinker's implementation. See
Section~\ref{sec:tinker_closed} for detail.

\paragraph{Cross-platform hardware confounding.}
Tinker, Modal, Colab, and the Institutional A100 node use different accelerators,
different memory hierarchies, and different inference/training stacks
(Tinker proprietary, Modal vLLM, Colab PyTorch-native, TRL+Accelerate).
Observed differences in reward dynamics are not purely algorithmic.

\paragraph{Carbon estimates are approximate.}
As discussed in Section~\ref{sec:carbon}, Tinker GPU-hour and TDP are
inferred, not measured. Grid intensity is region-average, not marginal.
Numbers in Table~\ref{tab:carbon2} should be treated as
order-of-magnitude.

\paragraph{Exploratory, not confirmatory.}
This is an exploratory case study, not a pre-registered confirmatory
experiment. We did not pre-register primary hypotheses. All $p$-values are
un-corrected for the number of comparisons actually made across the
paper; the Bonferroni-surviving comparisons in
Section~\ref{sec:results} are so flagged, but most of our secondary
analyses are descriptive.

\paragraph{Geographic and demographic limits.}
The specific city and country of the authors' institution are withheld
here for anonymity but disclosed in the non-anonymous version of the
paper. Our choice of benchmarks, prompt phrasings, and evaluation design
reflects that single-region context. Independent replication by groups
in other regions, on non-English tasks, and with non-Qwen / non-Llama
families is an open question.

 
\phantomsection\label{sec:acknowledgments}
\begin{ack}
Acknowledgments are withheld for blind review. We thank the Tinker team for
API access, the Modal team for GPU compute credits (H100 cluster), Weights
\& Biases for experiment tracking infrastructure, and Hugging Face for model
hosting. We are grateful to the open-source communities behind TRL, Stable
Baselines3, CleanRL, Tianshou, OpenRLHF, veRL, and the broader Hugging Face
ecosystem.
\end{ack}

\bibliographystyle{plainnat}
\bibliography{../references}

\appendix


\section{Run Manifest}
\label{sec:run-manifest}

This manifest summarises every (backend, model, task) cell in the benchmark
at submission time, organised by evidence tier. Cells marked \textbf{A} have at
least three completed seeds and support inference statistics. Cells marked
\textbf{B} are single-seed completed runs and are reported descriptively only;
where these single-seed runs are used as data points in within-task
correlation analyses (e.g.\ the ZVF $N{=}7$ stratified Pearson on
model-varying GSM8K), the cell carries a $\dagger$. Cells marked
\textbf{C} failed, partially completed, or never started; we list
them for honesty and exclude them from headline claims. ``Held-out'' indicates
that a held-out test-split evaluation exists for at least one seed in that cell.
``Logs'' indicates that step-level group reward traces are released.
``Ckpt'' indicates that a usable LoRA checkpoint is released.
``Reward'' is binary verifiable (BV), shaped (SH), or pairwise (PW). Raw
step-level reward traces are released for every Tier-A and Tier-B row.

\begin{table}[h]
\centering\scriptsize
\setlength{\tabcolsep}{3pt}
\begin{tabular}{lllllrrrcccc}
\toprule
Tier & Backend & Model & Task & Split & N & Seeds & Steps$_{\max}$ & Reward & Held-out & Logs & Ckpt \\
\midrule
\multicolumn{12}{l}{\textit{Tier A --- multi-seed completed cells, used for inference statistics}} \\
\midrule
A & Modal       & cleanrl-ppo               & math                & gen     & 5 & 5 & 97280  & BV & --- & \checkmark & --- \\
A & Modal       & sb3-ppo                   & math                & gen     & 5 & 5 & 100352 & BV & --- & \checkmark & --- \\
A & Modal       & tianshou-ppo              & math                & gen     & 5 & 5 & 100000 & BV & --- & \checkmark & --- \\
A & Modal L4    & qwen2.5-0.5b              & math                & gen     & 5 & 5 & 125    & BV & --- & \checkmark & \checkmark \\
A & Tinker      & qwen3-8b (Wave-6 ablation)& gsm8k               & train+held-out & 12 & 1 & 30 & BV & \checkmark & \checkmark & --- \\
A & Tinker      & llama-3.3-70b-inst        & gsm8k               & train   & 4 & 4 & 30     & BV & --- & \checkmark & --- \\
\midrule
\multicolumn{12}{l}{\textit{Tier B --- single-seed completed Tinker case studies}} \\
\midrule
B$^\dagger$ & Tinker      & qwen3-8b (G=8, base)      & gsm8k               & train+held-out & 1 & 1 & 30 & BV & \checkmark & \checkmark & --- \\
B$^\dagger$ & Tinker      & qwen3.5-4b                & gsm8k               & train   & 1 & 1 & 30 & BV & --- & \checkmark & --- \\
B$^\dagger$ & Tinker      & llama-8b-inst (gsm8k)     & gsm8k               & train   & 1 & 1 & 30 & BV & --- & \checkmark & --- \\
B$^\dagger$ & Tinker      & deepseek-v3.1             & gsm8k               & train   & 1 & 1 & 20 & BV & --- & \checkmark & --- \\
B$^\dagger$ & Tinker      & nemotron-120b             & gsm8k               & train   & 1 & 1 & 20 & BV & --- & \checkmark & --- \\
B$^\dagger$ & Tinker      & gpt-oss-20b               & gsm8k               & train   & 1 & 1 & 30 & BV & --- & \checkmark & --- \\
B$^\dagger$ & Tinker      & kimi-k2                   & gsm8k               & train   & 1 & 1 & 20 & BV & --- & \checkmark & --- \\
B & Modal H100  & llama-3.2-3b-inst         & gsm8k               & train   & 1 & 1 & 30 & BV & --- & \checkmark & --- \\
B & Modal H100  & llama-3.2-1b-inst         & gsm8k               & train   & 1 & 1 & 30 & BV & --- & \checkmark & --- \\
B & Modal H100  & qwen3-8b (TRL ref)        & gsm8k               & train   & 1 & 1 & 30 & BV & --- & \checkmark & --- \\
B & Modal H100  & llama-8b-inst             & gsm8k               & train   & 1 & 1 & 30 & BV & --- & \checkmark & --- \\
B & Tinker      & llama-8b-inst (tool)      & tool\_use           & train   & 1 & 1 & 30 & SH & --- & \checkmark & --- \\
B & Tinker      & qwen3-32b (tool)          & tool\_use           & train   & 1 & 1 & 30 & SH & --- & \checkmark & --- \\
\midrule
\multicolumn{12}{l}{\textit{Tier B --- Bitter Lesson Campaign (single-seed Tinker frontier runs, GSM8K, 30 steps, completed)}} \\
\midrule
B & Tinker      & qwen3-8b-base             & gsm8k               & train   & 1 & 1 & 30 & BV & --- & \checkmark & --- \\
B & Tinker      & kimi-k2-thinking          & gsm8k               & train   & 1 & 1 & 30 & BV & --- & \checkmark & --- \\
B & Tinker      & kimi-k2.5                 & gsm8k               & train   & 1 & 1 & 30 & BV & --- & \checkmark & --- \\
B & Tinker      & gpt-oss-120b              & gsm8k               & train   & 1 & 1 & 30 & BV & --- & \checkmark & --- \\
B & Tinker      & qwen3.5-397b              & gsm8k               & train   & 1 & 1 & 30 & BV & --- & \checkmark & --- \\
B & Tinker      & llama-3.1-70b             & gsm8k               & train   & 1 & 1 & 30 & BV & --- & \checkmark & --- \\
B & Tinker      & deepseek-v3.1-base        & gsm8k               & train   & 1 & 1 & 30 & BV & --- & \checkmark & --- \\
B & Tinker      & llama-3.1-8b              & gsm8k               & train   & 1 & 1 & 30 & BV & --- & \checkmark & --- \\
\midrule
\multicolumn{12}{l}{\textit{Tier B --- Group-Size Ablation (Qwen3-8B, GSM8K, single-seed per G, 30 steps, completed)}} \\
\midrule
B & Tinker      & qwen3-8b (G=2)            & gsm8k               & train   & 1 & 1 & 30 & BV & --- & \checkmark & --- \\
B & Tinker      & qwen3-8b (G=4)            & gsm8k               & train   & 1 & 1 & 30 & BV & --- & \checkmark & --- \\
B & Tinker      & qwen3-8b (G=8)            & gsm8k               & train   & 1 & 1 & 30 & BV & --- & \checkmark & --- \\
B & Tinker      & qwen3-8b (G=16)           & gsm8k               & train   & 1 & 1 & 30 & BV & --- & \checkmark & --- \\
B$^\dagger$ & Tinker      & qwen3-8b (G=32) & gsm8k               & train   & 1 & 1 & 30 & BV & --- & \checkmark & --- \\
C$^\S$ & ---         & qwen3-8b (G=64)$^\S$      & gsm8k (extrap.)     & ---     & 0 & 0 & 0  & --- & --- & --- & --- \\
\midrule
\multicolumn{12}{l}{\textit{Tier C --- failed, partial, or never-started cells (listed for honesty)}} \\
\midrule
C & External    & 0.5b-dpo                  & keyword\_generation & ---     & 1 & 0 & 0  & PW & --- & --- & --- \\
C & External    & 3b-tool                   & tool\_calling       & ---     & 1 & 0 & 0  & SH & --- & --- & --- \\
C & External    & qwen3-4b                  & math\_reasoning     & ---     & 1 & 0 & 0  & BV & --- & --- & --- \\
C & External    & qwen3-8b                  & code\_generation    & ---     & 1 & 0 & 0  & SH & --- & --- & --- \\
C & Modal H100  & llama-3.2-1b              & gsm8k (alt run)     & ---     & 1 & 0 & 0  & BV & --- & --- & --- \\
C & Modal H100  & llama-3.2-3b              & gsm8k (alt run)     & ---     & 1 & 0 & 0  & BV & --- & --- & --- \\
C & Modal H100  & qwen3-8b                  & humaneval           & ---     & 1 & 1 & 0  & BV & --- & --- & --- \\
C & Modal H100  & qwen3-8b                  & kl\_tracking        & ---     & 1 & 1 & 0  & ---& --- & --- & --- \\
C & Tinker      & qwen3-235b-moe            & gsm8k (partial)     & train   & 1 & 1 & 15 & BV & --- & \checkmark & --- \\
C & Tinker      & qwen3-30b-moe             & gsm8k (partial)     & train   & 1 & 1 & 20 & BV & --- & \checkmark & --- \\
C & Tinker      & qwen3-30b-moe-inst        & gsm8k (partial)     & train   & 1 & 1 & 10 & BV & --- & \checkmark & --- \\
C & Tinker      & qwen3-32b                 & gsm8k (partial)     & train   & 1 & 1 & 10 & BV & --- & \checkmark & --- \\
C & Tinker      & qwen3.5-27b               & gsm8k (partial)     & train   & 1 & 1 & 10 & BV & --- & \checkmark & --- \\
\bottomrule
\end{tabular}
\caption{Per-(backend, model, task) run manifest at submission time.
Tier classification: \textbf{A}~$\geq3$ seeds completed (or, for the
Wave-6 single-seed sensitivity ablation, $\geq12$ runs at fixed seed
sweeping a hyperparameter), used for inference; \textbf{B}~single seed
completed, used descriptively; \textbf{C}~failed, partial, or not
started, listed for honesty. Rows marked $\dagger$ in the Tinker case-study
block are the seven single-seed Tinker GSM8K cells used as data points in
the ZVF model-varying within-task stratified Pearson ($N{=}7$,
Section~\ref{sec:zvf-analysis} of the ZVF variant); the additional
$\dagger$ on the G=32 group-size row marks it as the single-seed cell
that anchors the upper end of the within-sweep token-budget reanalysis
in Appendix~\ref{sec:appendix:group-size-reconcile}.
``Split'' is the evaluation distribution used for reported numbers
(\emph{train}: training-prompt reward; \emph{train+held-out}: also
evaluated on a held-out test split; \emph{gen}: synthetic generated
data; ``---'' for failed cells). ``Reward'' is BV (binary verifiable),
SH (shaped), or PW (pairwise). All Tinker rows depend on a closed-source
backend. The Bitter Lesson Campaign and Group-Size Ablation rows are
the single-seed Tinker frontier and group-size evidence underlying
Table~4 of the DNB variant; the Tier-A Modal L4 Qwen2.5-0.5B row is
the workshop-headline multi-seed reference run.\\[2pt]
\textit{$^\S$~G=64: no direct training run was executed at this group size.
Numbers appearing in the G=64 column of the token-budget reanalysis
(Appendix~\ref{sec:appendix:group-size-reconcile},
Table~\ref{tab:groupsize-tokennorm}) are extrapolated from the
inverted-U fit on G$\in\{4,8,16,32\}$ and are descriptive only; we
explicitly do \emph{not} treat G=64 as Tier-B evidence.}}
\label{tab:run-manifest}
\end{table}
 \section{Formal Definition of the Zero-Variance Fraction, Non-Tautology, and Scope}
\label{sec:appendix:zvf-formalization}

\subsection{Formal Definition}

Let a GRPO batch $B_t$ at optimizer step $t$ consist of $|B_t|$ prompt-groups,
each group $g$ containing $K$ rollouts with scalar outcome rewards
$r_{g,1}, \ldots, r_{g,K} \in \mathbb{R}$. We define the \emph{Zero-Variance
Fraction} (ZVF) as
\begin{equation}
\label{eq:zvf}
\mathrm{ZVF}_t \;=\; \frac{1}{|B_t|} \sum_{g \in B_t}
  \mathbb{1}\!\left[\;\widehat{\mathrm{Var}}(r_{g,1}, \ldots, r_{g,K}) \le \varepsilon\;\right],
\end{equation}
where $\widehat{\mathrm{Var}}$ is the unbiased sample variance and
$\varepsilon$ is a small numerical tolerance
(we use $\varepsilon = 10^{-6}$ for rewards normalized to $[0,1]$).
Intuitively, $\mathrm{ZVF}_t$ is the fraction of groups at step $t$ whose
rollouts all collapsed to the same outcome; those groups contribute zero
group-relative advantage and thus zero gradient signal to the GRPO update.

\subsection{Why ZVF Is Not a Tautological Re-expression of Mean Reward}

A reviewer may reasonably suspect that under binary outcome rewards
$r \in \{0,1\}$ and group-relative advantages, $\mathrm{ZVF}_t$ is a direct
function of the batch mean reward $\bar{r}_t$: if every group succeeds
(or fails) uniformly, $\mathrm{ZVF}_t = 1$. This is true at the extremes
$\bar{r}_t \in \{0,1\}$ but does \emph{not} hold at intermediate values.
For $K$ i.i.d.~Bernoulli rollouts with prompt-specific success probability
$p_g$, the expected ZVF under a null independence model is
\begin{equation}
\mathbb{E}[\mathrm{ZVF}_t] \;=\; \mathbb{E}_{p_g}\!\left[\,p_g^{K} + (1-p_g)^{K}\,\right].
\end{equation}
The right-hand side depends on the full prompt-difficulty distribution
$\{p_g\}$, not only on its first moment. Two populations can therefore
share the same mean reward $\bar r = 0.5$ yet have very different ZVF:
a bimodal population with $p_g \in \{0,1\}$ yields $\mathrm{ZVF}=1$,
while a uniform-hard population with all $p_g = 0.5$ and $K=8$ yields
$\mathrm{ZVF} \approx 2 \cdot (0.5)^8 \approx 0.008$. The substantive
claim is thus not ``high mean reward implies high ZVF,'' but rather that
ZVF is sensitive to \emph{how} success mass is distributed across prompts.

\subsection{What the Released Artifact Establishes}

The released artifact bundled with this paper contains aggregate reward
trajectories, held-out checkpoint summaries, and the cross-framework parser
specification released with the anonymous artifact. It does
\emph{not} contain the per-group step-level rollout logs needed to support
a new partial-correlation table that controls simultaneously for
batch-mean reward, entropy, KL drift, and advantage variance. We therefore
do \emph{not} claim a fresh incremental-$R^2$ estimate here, and we mark
any implication that the current artifact proves a new partial-correlation
result beyond the already released summary traces as rerun required.

The narrower empirical claim made in the revised paper is the one directly
supported by the released evidence: in outcome-reward GRPO, high ZVF
coincides with the disappearance of within-group contrast, and at any step
where $\mathrm{ZVF}=1$ the group-relative advantage is mechanically zero so
no gradient flows from that batch. We additionally observe co-occurrence
between elevated ZVF and reward plateaus or regressions in some runs, but
we do not establish a causal link: the within-task stratified analysis in
Section~\ref{sec:zvf-analysis} does not support a within-task predictive
relationship between mean ZVF and last-10 reward. The operational meaning
of ZVF in the present benchmark is therefore the local mechanical claim,
not a performance-causal one.

\subsection{Boundary Conditions and Scope}

We state explicitly when ZVF is \emph{not} informative, partly in response
to the tautology concern:
\begin{itemize}[leftmargin=1.2em,itemsep=0pt]
  \item \textbf{Dense or continuous rewards.} When rewards are continuous
        (e.g.,~process-reward models or graded code-execution partial credit),
        $\widehat{\mathrm{Var}}$ is almost surely nonzero and
        $\mathrm{ZVF} \to 0$. ZVF degenerates and must be replaced by a
        per-step-reward-variance diagnostic or alternative diagnostic surrogates.
  \item \textbf{$K=1$.} Group variance is undefined and ZVF is trivially $1$.
  \item \textbf{SFT-saturated baselines.} When the policy already solves
        the task ($\bar p_g \to 1$ for essentially every prompt),
        $\mathrm{ZVF} \to 1$ and the diagnostic loses discriminative power;
        in this regime RL no longer provides useful signal anyway.
  \item \textbf{Non-outcome-reward RL.} For DPO- or regression-style
        training there are no rollout groups and ZVF is ill-defined.
\end{itemize}
Within its intended scope -- outcome-reward GRPO on verifiable tasks with
$K \ge 4$ and non-saturated reward support -- ZVF is a cheap, model-agnostic
early-warning diagnostic. Outside that scope, the paper now recommends
explicit surrogates instead of over-claiming portability.
 \section{Reconciling the Group-Size Result: Token-Budget-Normalized Sweeps}
\label{sec:appendix:group-size-reconcile}

\subsection{The Apparent Contradiction}

A reviewer correctly identified that Table~6 reports $G{=}4$ as the
highest last-10 reward ($52.1\%$) under a fixed-step budget, while the main text
claims that $G{\approx}32$ ``maximizes gradient utilization''. These
observations are compatible but require the axis of comparison to be
made explicit. Under a \emph{fixed number of optimizer steps}, larger
$G$ costs proportionally more tokens per step, so a step-budget
comparison benefits small $G$; under a \emph{fixed total token budget},
larger $G$ amortizes wasted (all-correct or all-incorrect) rollouts
and yields higher effective gradient signal per token.

\subsection{Formal Definition of Token-Budget Gradient Efficiency}

To analyze the token-budget tradeoff we define a \emph{token-budget gradient
efficiency} $\mathrm{GE}_{\text{tok}}$---a per-token quantity distinct from the
main-text gradient utilization $\mathrm{GU}{=}1{-}\mathrm{ZVF}$ (a unit-free
fraction); we use a separate symbol to avoid conflating the two. It is the
expected squared policy-gradient-norm contribution per unit of token budget,
restricted to rollouts whose group advantage is nonzero:
\begin{equation}
\label{eq:gu}
\mathrm{GE}_{\text{tok}}(G, B) \;=\;
\frac{\displaystyle\mathbb{E}\!\left[\,\|\nabla_\theta \mathcal{L}\|_2^{2}\,\cdot\,
  \mathbb{1}\!\left[A_g \ne 0\right]\,\right]}
{\displaystyle G \cdot K \cdot \bar{L}_\text{rollout}},
\end{equation}
where $A_g$ is the within-group advantage, $K$ the rollouts per group,
and $\bar{L}_\text{rollout}$ the mean rollout length.
Intuitively, $\mathrm{GE}_{\text{tok}}$ penalizes compute spent on rollouts that
contribute zero gradient ($\mathbb{1}[A_g\!=\!0]$, i.e.~ZVF-group
members) and rewards concentrating token budget where advantage is
informative. An estimator that avoids explicit gradient storage uses
advantage variance as a proxy:
$\widehat{\mathrm{GE}}_{\text{tok}}(G,B) \!\propto\!
 (1 - \mathrm{ZVF}) \cdot \widehat{\mathrm{Var}}_A / (G\cdot K \cdot \bar{L}_\text{rollout})$.
Because the $G$ in the denominator grows faster than the numerator saturates,
$\widehat{\mathrm{GE}}_{\text{tok}}$ is monotone decreasing in $G$ in our sweep.

\subsection{Token-Budget-Normalized Sweep}

We re-analyze the $G$-sweep under a fixed \emph{token} rather than
\emph{step} budget. Let $T$ denote total optimizer-visible tokens.
At fixed $T$, the number of optimizer steps is
$n_\text{step}(G) = T / (G\cdot K\cdot \bar{L})$, which decreases
with $G$. The question is whether the larger per-step gradient
signal of larger $G$ outweighs the reduced number of updates.

\begin{table}[ht]
\centering
\small
\begin{tabular}{lcccccc}
\toprule
Total tokens $T$ & Metric &
$G{=}4$ & $G{=}8$ & $G{=}16$ & $G{=}32$ & $G{=}64$ \\
\midrule
\multirow{2}{*}{$1\text{M}$}  & heldout acc. & $0.41$ & $\mathbf{0.48}$ & $0.47$ & $0.42$ & $0.35$ \\
                               & $\widehat{\mathrm{GE}}_{\text{tok}}$ ($\times 10^{-3}$) & $\mathbf{1.72}$ & $1.12$ & $0.66$ & $0.34$ & $0.17$ \\
\multirow{2}{*}{$4\text{M}$}  & heldout acc. & $0.55$ & $0.67$ & $\mathbf{0.69}$ & $0.66$ & $0.58$ \\
                               & $\widehat{\mathrm{GE}}_{\text{tok}}$ ($\times 10^{-3}$) & $\mathbf{2.13}$ & $1.46$ & $0.87$ & $0.48$ & $0.25$ \\
\multirow{2}{*}{$16\text{M}$} & heldout acc. & $0.63$ & $0.78$ & $0.82$ & $\mathbf{0.84}$ & $0.80$ \\
                               & $\widehat{\mathrm{GE}}_{\text{tok}}$ ($\times 10^{-3}$) & $\mathbf{2.33}$ & $1.61$ & $0.96$ & $0.56$ & $0.29$ \\
\multirow{2}{*}{$64\text{M}$} & heldout acc. & $0.64$ & $0.80$ & $0.85$ & $\mathbf{0.88}$ & $0.87$ \\
                               & $\widehat{\mathrm{GE}}_{\text{tok}}$ ($\times 10^{-3}$) & $\mathbf{2.42}$ & $1.66$ & $0.99$ & $0.58$ & $0.32$ \\
\bottomrule
\end{tabular}
\caption{Token-budget-normalized $G$-sweep for GRPO on Qwen3-8B / GSM8K
($K{=}1$ rollout per sample, $\bar L = 384$). \emph{Illustrative reanalysis:}
the held-out accuracy cells are reconstructed point estimates from existing GRPO
ablation logs (the \texttt{FALLBACK\_ROWS} table in the script), \emph{not}
freshly measured per-seed runs; per-cell seed counts are not retained and no
measured $G\times T$ sweep over $\{4,8,16,32,64\}$ was executed. Bold in the
accuracy row marks the per-row accuracy maximum, which shifts rightward with
$T$: $G{=}8$ at $T{=}1\text{M}$ (near the fixed-\emph{step} optimum $G{=}4$ of
Table~6), $G{=}16$ at $4\text{M}$, and
$G{\approx}32$ at $T{\ge}16\text{M}$ (the canonical training budget) \emph{in
the reconstruction}. This is a hypothesis for a measured sweep, not an
empirical optimum. The
token-budget gradient-efficiency estimator $\widehat{\mathrm{GE}}_{\text{tok}}$
(Eq.~\ref{eq:gu}; distinct from the main-text $\mathrm{GU}{=}1{-}\mathrm{ZVF}$)
is by contrast \emph{monotone decreasing} in $G$ at every budget (bold marks its
maximum, always $G{=}4$): per-token gradient efficiency favors small $G$ and does
\emph{not} on its own justify large $G$. See
\texttt{experiments/group\_size\_token\_normalized.py}.}
\label{tab:groupsize-tokennorm}
\end{table}

\paragraph{Accuracy-optimal $G$ shifts with budget.}
The per-row accuracy-maximizing group size increases with the token budget
($G^\star = 8, 16, 32, 32$ at $T = 1, 4, 16, 64\,\text{M}$) in the reconstructed
grid, so the hypothesis is budget-dependent. We report this as a qualitative
argmax pattern to test---no trained matched-budget $G{=}32$ arm validates it and
no committed script fits a parametric
inverted-U envelope or its bootstrap confidence interval, so we make no such
quantitative claim.

\subsection{Reconciliation Statement}

We therefore revise the main-text claim as follows.
\begin{quote}\itshape
Under a fixed \emph{step} budget at small total token counts,
$G{=}4$ attains the highest last-10 reward ($52.1\%$, Table~6).
Under fixed \emph{total tokens} at the canonical training scale
used elsewhere in the paper ($T\!\ge\!16$M), $G{\approx}32$
maximizes reconstructed held-out accuracy in this illustrative reanalysis; it
is not a measured recommendation. (The
per-token gradient-efficiency estimator $\widehat{\mathrm{GE}}_{\text{tok}}$ is
\emph{not} co-maximized: it decreases monotonically in $G$, so it favors small
$G$ and is reported only to show that the accuracy gain at large $G$ comes at a
per-token efficiency cost.) The accuracy-optimal $G$ shifts rightward with $T$,
so the recommended $G$ depends on the practitioner's compute budget, not on a
universal heuristic.
\end{quote}
The reviewer's concern is taken seriously: the ``more rollouts is
always better'' heuristic is false, \emph{and} the opposite heuristic
``small $G$ is always better'' is equally false; the correct claim is
that there exists a budget-dependent optimum and practitioners should
locate it via the reanalysis script
\texttt{experiments/group\_size\_token\_normalized.py}.

\subsection{Measured Small-Scale Confirmation of the Inverted-U}
\label{sec:appendix:group-size-measured}

Because the token-normalized table above is an illustrative reanalysis, we
separately ran a \emph{measured} group-size sweep to confirm the qualitative
inverted-U with fresh per-seed data (Table~\ref{tab:groupsize-measured};
driver \texttt{experiments/modal/modal\_groupsize\_zvf\_sweep.py}, artifact
\texttt{experiments/results/groupsize\_zvf\_sweep.tsv}). On the ungated
\texttt{Qwen2.5-0.5B} with a verifiable correctness reward on two-operand
addition, $G \in \{2,4,8,16\}$ at three seeds each, held-out accuracy traces an
inverted-U shape with its numerical maximum at the intermediate $G{=}8$
($0.982 \to 0.988 \to 0.990 \to 0.978$); we emphasize the apex is \emph{not}
statistically separable from its neighbours ($G{=}4$ $0.988{\pm}.002$ vs.\
$G{=}8$ $0.990{\pm}.003$, overlapping SEs at $n{=}3$), so the supported claim is
only that the optimum is interior, not that it is precisely $G{=}8$. The mean
zero-variance fraction falls monotonically with $G$ ($0.84 \to 0.76 \to 0.69 \to
0.63$), the direction $\mathrm{ZVF}(p,G)=p^{G}+(1-p)^{G}$ predicts (the closed
form captures the trend, not the absolute level---see the formalization
appendix). This is a small model on an easy task (accuracies are near ceiling,
so the effect sizes are modest) and is \emph{not} a substitute for a
Qwen3-8B/GSM8K sweep; we report it only as a measured, reproducible indication
that the optimal $G$ is interior rather than smallest-or-largest.

\begin{table}[ht]
\centering
\small
\begin{tabular}{lcccc}
\toprule
Metric & $G{=}2$ & $G{=}4$ & $G{=}8$ & $G{=}16$ \\
\midrule
Held-out acc.\ (mean $\pm$ SE, $n{=}3$) & $0.982{\pm}.004$ & $0.988{\pm}.002$ & $\mathbf{0.990}{\pm}.003$ & $0.978{\pm}.006$ \\
Mean ZVF & $0.84$ & $0.76$ & $0.69$ & $0.63$ \\
\bottomrule
\end{tabular}
\caption{Measured group-size sweep on \texttt{Qwen2.5-0.5B} / arithmetic
(correctness reward, 40 steps, 3 seeds per $G$). Held-out accuracy peaks at the
intermediate $G{=}8$; ZVF decreases monotonically with $G$. Small-scale,
ungated confirmation of the inverted-U---not a GSM8K-scale claim.}
\label{tab:groupsize-measured}
\end{table}
 
\section*{Stratified Diagnostics}
\label{sec:stratified-diagnostics}

\paragraph{Bootstrap CIs and permutation tests.}
Pearson $r$ between mean ZVF and last-10 reward on each of the four cell-pools
used in this paper, with 10\,000-resample bootstrap 95\% CIs and
10\,000-permutation $p$-values (Table~\ref{tab:zvf-bootstrap}). The pooled
cross-task signal is large; the within-task subsets cross zero.

\begin{table}[h]
\centering\small
\begin{tabular}{lrrrr}
\toprule
Analysis & $r$ & $N$ & 95\% CI (bootstrap) & $p_{\text{perm}}$ \\
\midrule
Pooled original                 & $-0.769$ & $15$ & $[-0.93,\,-0.52]$ & $0.0004$ \\
Pooled extended                  & $-0.443$ & $29$ & $[-0.69,\,+0.03]$ & $0.0146$ \\
Within GSM8K                    & $+0.401$ & $19$ & $[-0.08,\,+0.83]$ & $0.0837$ \\
GSM8K model-varying\textsuperscript{\dag}    & $+0.126$ & $7$  & $[-0.97,\,+1.00]$ & $0.8139$ \\
\bottomrule
\end{tabular}
\caption{ZVF--last-10 reward correlations with bootstrap CIs and permutation
$p$-values. The extended pool reports every cell whose step-level ZVF trace is
in the released artifact. \textsuperscript{\dag}~The $N{=}7$ nonparametric CI
is uninformative; the parametric Fisher-$z$ CI is $[-0.69,\,+0.80]$, also
crossing zero.}
\label{tab:zvf-bootstrap}
\end{table}

\paragraph{Partial correlations.}
With task fixed effects, log group size, and pre-RL accuracy as controls
(Table~\ref{tab:zvf-partials}); we never control for last-10 reward, the
outcome. Adding task FE alone collapses the pooled $r$ from $-0.44$ to
$-0.03$; with any combination of $\log G$ and base accuracy the partial $r$
is between $+0.02$ and $+0.14$ and never significant.

\begin{table}[h]
\centering\small
\begin{tabular}{lrrrl}
\toprule
Specification & $r_{\text{partial}}$ & df & $p$ & Interpretation \\
\midrule
Raw pooled, no controls                          & $-0.443$ & $27$ & $0.016$ & Task is a lurking variable. \\
$+$~task FE                                      & $-0.026$ & $26$ & $0.896$ & Task identity removes the sign. \\
GSM8K within-task                                & $+0.401$ & $17$ & $0.088$ & Sign flips, CI crosses zero. \\
GSM8K model-varying                              & $+0.126$ & $5$  & $0.788$ & Underpowered ($N{=}7$). \\
Pooled $+$ task FE $+$ $\log G$                  & $+0.021$ & $25$ & $0.916$ & Group size does not rescue it. \\
Pooled $+$ task FE $+$ base acc                  & $+0.135$ & $25$ & $0.502$ & Pre-RL ability does not rescue it. \\
Pooled $+$ task FE $+$ $\log G$ $+$ base acc     & $+0.136$ & $24$ & $0.507$ & Joint controls; not significant. \\
GSM8K $+$ $\log G$ $+$ base acc                  & $+0.420$ & $15$ & $0.094$ & Within-task partial agrees with raw. \\
\bottomrule
\end{tabular}
\caption{Partial Pearson correlations of mean ZVF with last-10 reward
(\texttt{pingouin.partial\_corr}). Last-10 reward is the outcome and is never
on the right-hand side. Stack fixed effects are not reported (too little
within-stack variation: $\le 3$ stacks with $\ge 2$ runs).}
\label{tab:zvf-partials}
\end{table}

\paragraph{Interpretation.}
The pooled correlation does not survive stratification: with task fixed
effects $r$ collapses to $-0.03$, and on the largest within-task subset the
sign flips to $+0.40$. The remaining within-GSM8K signal is underpowered,
with a 95\% CI that crosses zero. This is not evidence of absence---it is
evidence that the pooled monotone story is not supported once the task
confound is removed.
 
\end{document}
